\chapter{Conclusion and Future Work}
\label{chap:conclusion-future-work}

\section{Conclusion}
\label{sec:conclusion}

This project investigated how prism adaptation can be implemented and studied in VR when the field does not use one shared terminology or one shared procedural framework. The literature review showed that studies often use overlapping terms for different implementations, including full-field visual shifts, effector-feedback shifts, virtual hand shifts, and visuomotor rotations. Based on this gap, the project developed a configurable VR prototype that made these implementation choices explicit rather than hiding them behind a single ``prism effect'' label.

The final prototype implemented three prism-like perturbations, Translation, Rotation, and Skew, together with two assessment tasks, Open-Loop Pointing and Line Bisection, and an exposure task for repeated visuomotor correction. The system also supported both controller input and embodied hand tracking, logged participant sessions to CSV, and provided an RShiny dashboard for summary, comparison, participant-level, quality, spatial, and trajectory analysis. In this sense, the project achieved its practical goal: it produced a working sandbox for comparing VR prism-effect implementations through a shared logging and analysis pipeline.

The participant evaluation showed that the system could collect complete data across 10 participants and 80 finished modules. The data contained measurable baseline-to-post shifts, especially in the controller conditions where all perturbations were above the participant-specific no-effect reference at the median level. This supports the prototype's main purpose: it can run multiple VR prism-effect implementations inside one shared workflow and produce comparable prism-adaptation-like outcome data. The evaluation also showed why the framework needs detailed logging and quality views. The None condition produced non-zero shifts, and several participant-specific factors affected the results, including rushed exposure, fatigue, body compensation, hand posture, tracking instability, and possible carry-over between modules. The SSQ responses showed a small overall increase in reported discomfort after testing, but simulator sickness was not the main limitation observed during the evaluation.

The broader contribution is methodological. The project shows that VR prism adaptation should be described through explicit dimensions such as perturbation class, exposure visibility, embodiment representation, action-confirmation method, task family, reset procedure, and participant posture. Without these distinctions, different implementations risk being treated as equivalent even when they produce different participant behaviour and different analysis outcomes.

\section{Future Work}
\label{sec:future-work}

Future work should first improve the evaluation procedure. The next study should use stricter instructions, longer familiarization, clearer exclusion or repetition criteria, and stronger control of posture and pace. Participants should be prevented from rushing through exposure, and modules with very low hit rates or clear calibration mistakes should either be repeated immediately or marked separately before analysis. The procedure should also include stronger reset or washout logic between perturbations, since the current results suggested possible carry-over between modules. SSQ should remain part of the procedure, but it should be reported alongside a separate fatigue or workload measure so that cybersickness can be distinguished from raised-arm tiredness.

The second priority is to refine the interaction design. The current pointer-ray implementation made the tasks feasible, but it also allowed participants to aim through wrist, arm, or torso compensation. A future version should either constrain the task closer to a forward reaching movement or log the relevant body segments so that compensation can be detected. Embodied hand tracking should also be made more robust through posture guidance, tracking-quality warnings, and clearer calibration feedback. The audio confirmation feedback was well received and should be retained.

The third priority is to expand the prototype as a framework tool. Exposure visibility should become a configurable first-class variable, making it possible to compare concurrent feedback against terminal feedback in the same system. Miss handling should also be configurable, so researchers can choose between counting every attempt, repeating targets until hit, or using a hybrid rule. The analysis dashboard should continue to emphasize participant-level quality and no-effect drift, because these were essential for interpreting the noisy data.

Finally, future evaluations should involve a larger participant sample and stronger separation between conditions. The current study was appropriate for testing the prototype with healthy participants, but it was still a single-session evaluation with several perturbations performed back-to-back. A later rehabilitation-oriented study could test whether the same framework supports patient-facing procedures, subtype-sensitive assessment, and repeated-session protocols. Before that step, the sandbox should be stabilized further so that procedural noise does not obscure the effect of the perturbation itself.
